% Pursuing Stacks (cotes 134-2 to 134-8): archive cote 134-5, site page 43 (PDF page 44), typescript page 230.
% typescript. Draft compared: ps4.tex:2089-2156.
% AI-transcribed by Claude Opus 5.5 (claude-opus-5-5, adaptive thinking, effort high), Message Batches API, batch ps2 (msgbatch_01NVwr1Yuo8qJf5rykeGBXPK), 2026-11-12; not yet checked by eye.
\scanpage{134-5/43}
\mspage{230}

(Compare with propositions on pages 121 and 143\hand{.}) Whereas this \inserted{latter} assumption
admittedly is quite a stringent one, it is however of a type which has
become familiar to us in relation with test categories, where it seems a
rather common lot.

If B satisfies these conditions (as in prop.4) and if A is totally
aspheric, we see from prop.4 that the category $\underline{\mathrm{Asph}}(A,B)$ of aspheric
functors from A to B is stable under binary products. Now \xout{itxisxeasily}
\inserted[typed]{a non-empty object C of (Cat)}
\xout{\illegible{4}} stable under binary products gives rise to a
category $C\hat{}$ which is clearly totally aspheric for any basic localizer $\underline{W}$ ,
and in particular for the usual one \unsure{$\underline{W}_e$|$\underline{W}_0$} corresponding to usual weak equi-%
valence. A fortiori such a category C \inserted[typed]{\xout{(ifxnonxempty)}} is 1-connected \hand{$\leftarrow$} \hand{(}which is easily
checked \inserted{too} by down-to-earth direct arguments\hand{)} \deleted{too}. Thus the reflections of
page 225 apply, and imply that if $\underline{\mathrm{Asph}}(A,B)$ is non-empty i.e. if there
is at least \uline{one} aspheric functor $i\colon A\longrightarrow B$, then there is a canonical
trans\hand{i}tive system of \unsure{\over{si}{is}}omorphisms between all functors
\[ (\ast) \qquad \mathrm{Hot}_B \longrightarrow \mathrm{Hot}_A \]
\margin{$\overline{i^*}$} of the type \over{$i^*$}{$\overline{i^*}$}, and hence there is a canonical functor $(\ast)$, defined up
to unique isomorphism, as announced in the remark on page 225.
\margin{For a generalization, see section 82; p. 272.}

13.6.

75. The generalities on aspheric maps of the last three sections should
be more than what is needed to develop now the notion of \uline{asphericity
structure} - which, \unsure{\hand{,}}together with the \inserted[typed]{closely related} notion of \uline{contractibility struc-%
ture}\hand{,} tentatively dealt with before, and the various \inserted[typed]{``test-}notions''\hand{(}\unsure{\over{cf}{e.g}}. \uline{test}
\deleted{functors} \uline{categories} and \uline{test\xout{x}functors}) seem\deleted{s} to me \deleted{at present} the main
pay off \deleted{till now} \inserted{so far} of our effort to come to a grasp of a general formalism
of ``homotopy models''.

\inserted{In the case of asphericity structures, just} \hand{(}\unsure{\over{A}{a}}s for the kindred notion of a contractibility structure, in all \deleted{cases} \inserted{instances}
I could think of at present, the underlying category M of an asphericity
structure is \uline{not} an object in (Cat) nor even a ``small category'', but is
``large'' - namely the cardinality of Ob M and F\unsure{l}(M) \unsure{\over{is}{are}} not in the\hand{``}univer-%
se''\inserted{$\underline{U}$} we are working in, \xout{whichxhasxxx} still less is M an object of $\underline{U}$ -
all we need\inserted{instead}, as usual, is that M be a \uline{$\underline{U}$-category}, namely that for any
two objects of $\underline{U}$ , $\underline{\mathrm{Hom}}(x,y)$ be an element of $\underline{U}$ . Till now, the universe
$\underline{U}$ has been present in our reflections in a very much implicit way, in kee-%
ping with the unformal nature of the reflections, which however by and by
have become more formal\deleted{,} \hand{(}as I finally let myself become invo\over{v}{l}ved in
\margin{43}
